\scanpage{33}
\mspage{14}

\uline{Corollaire} Supposons que $M$ soit \deleted{de la} isomorphe, soit au monoïde $\mathrm{Ep}(E)$ où $E$ est un ens infini, soit au monoïde opposé à $\mathrm{Mon}(E)$, avec $E$ un ens. infini. Alors $M$ est pseudocofilt\deleted{\unsure{rant}}\inserted{(il admet $=$ des éléments minimaux)} mais non cofiltrant (en fait $\unsure{f|g}uf=vf\Longrightarrow u=v$), $Y=B_M$ est $1$-connexe, mais il existe $F\in\mathrm{Ob}\,Y^\wedge$ tel que $X=Y_{/F}$ soit $0$-connexe et non $1$-connexe.

Pour ceci, il suffit de prendre \deleted{pour} la construction de la proposition, en prenant \deleted{\unsure{pour}} $\psi=\varphi$ n'importe quel objet minimal de $M$.

Il suffit de vérifier qu'on ne peut avoir $\varphi^{\deleted{\unsure{p}}r}=\varphi^s$ avec $r\ne s$. Mais si $r=s+h$, $h>0$, on aurait $\varphi^h\varphi^s=\mathrm{id}\,\varphi^s$, d'où $\varphi^h=1$, donc $\varphi$ serait inversible.
% letterhead
\par\smallskip
{\footnotesize\noindent\begin{minipage}[t]{.46\linewidth}\raggedright
\deleted{MINISTERE DE L'EDUCATION NATIONALE}\\ \deleted{ACADEMIE D'AIX-MARSEILLE}\\ \rule{2.6em}{.3pt}\\ \deleted{Inspection Académique}\\ \deleted{de Vaucluse}\\ \rule{2.6em}{.3pt}\\ \deleted{Tél. : 90 82 99 76}\\ \deleted{Référence : Secrétariat n\textsuperscript{o}}
\end{minipage}\hfill%
\begin{minipage}[t]{.51\linewidth}\raggedright
\deleted{AVIGNON, le}\\ \deleted{L'Inspecteur d'académie,}\\ \deleted{directeur des services départementaux}\\ \deleted{de l'Education Nationale de Vaucluse,}\\ \deleted{à M}
\end{minipage}\par}

% archive sheet number: 32
